\section{Součin rozpustnosti}

\subsection{Rozpustnost}
Udává kolik gramů látky se za dané teploty rozpustí ve 100 g rozpouštědla, zpravidla vody. Rozpustnost se zvyšující se teplotou zpravidla roste.

\begin{tabular}{|l|l|l|l|}
	\hline
	& \ce{NH4Cl} & NaCl & KCl \\\hline
	0 $^\circ$C & 29,41 & 35,63 & 28,14 \\\hline
	50 $^\circ$C & 50,42 & 36,70 & 43,05 \\\hline
	60 $^\circ$C & 55,30 & 37,08 & 45,88 \\\hline
\end{tabular}

\begin{itemize}
	\item \textbf{Nenasycený roztok} -- koncentrace látky je nižší, než její rozpustnost za dané teploty.
	\item \textbf{Nasycený roztok} -- koncentrace látky je shodná s její rozpustností za dané teploty.
	\item \textbf{Přesycený roztok} -- koncentrace látky je vyšší, než její rozpustnost za dané teploty.
\end{itemize}

\newpage

\subsection{Součin rozpustnosti}
Popisuje koncentrace iontů v nasycených roztocích.

\ce{AgNO3 + KCl -> AgCl v + KNO3}\\
\ce{AgCl <=> Ag+ + Cl-}

Rovnováhu v nasyceném roztoku chloridu stříbrného můžeme popsat pomocí rovnovážné konstanty:

$K = \frac{[\ce{Ag+}][\ce{Cl-}]}{[\textrm{AgCl}]}$

Z rovnovážné konstanty pak odvodíme \textit{součin rozpustnosti}:

$K_S = K.[\textrm{AgCl}] = [\ce{Ag+}][\ce{Cl-}]\\
pK_S = -\log K_S$

Hodnoty součinů rozpustnosti najdeme v chemických tabulkách.

\begin{center}
\begin{tabular}{|l|r@{,}l||l|r@{,}l|}
	\hline
	\textbf{Látka} & \multicolumn{2}{|c||}{\textbf{pK$_S$}} & \textbf{Látka} & \multicolumn{2}{|c|}{\textbf{pK$_S$}} \\\hline
	AgBr & 12 & 31 & CdS & 26 & 10 \\\hline
	AgCl & 9 & 75 & \ce{CoCO3} & 12 & 84 \\\hline
	AgI & 16 & 08 & CuBr & 8 & 28 \\\hline
	AgSCN & 11 & 97 & CuS & 35 & 20 \\\hline
	\ce{Ag2CrO4} & 11 & 61 & \ce{Ni(OH)2} & 15 & 50 \\\hline
	\ce{Ag2S} & 49 & 20 & PbS & 26 & 60 \\\hline
	\ce{Ag2SO4} & 4 & 77 & \ce{Sn(OH)4} & 56 & 00 \\\hline
	\ce{BaCO3} & 8 & 29 & \ce{Tl2S} & 20 & 30 \\\hline
	\ce{BaCrO4} & 9 & 93 & \ce{Zn3(AsO4)2} & 27 & 89 \\\hline
\end{tabular}
\end{center}

\subsubsection{Koncentrace iontů}

\zadani{Vypočítejte koncentraci iontů v nasyceném roztoku AgCl}

Chlorid stříbrný v roztoku disociuje podle rovnice:

\ce{AgCl <=> Ag+ + Cl-}

Součin rozpustnosti vypočítáme:

$K_S = 10^{-9,75} = 1,78\times 10^{-10}\\
K_S = [\ce{Ag+}][\ce{Cl-}] = x^2\\
x = \sqrt{K_S} = \sqrt{1,78\times 10^{-10}} = 1,33\times 10^{-5}$ M

\odpoved{Koncentrace stříbrného a chloridového iontu jsou 1,33$\times$10$^{-5}$ mol.dm$^{-3}$}

\clearpage